\documentclass[11pt]{article}

\usepackage[a4paper,margin=30mm]{geometry}
\usepackage{amsmath,amssymb,amsthm,mathtools}
\usepackage{booktabs}
\usepackage{enumitem}
\usepackage{xcolor}
\usepackage[hidelinks]{hyperref}
\usepackage{microtype}

\newtheorem{theorem}{Theorem}[section]
\newtheorem{lemma}[theorem]{Lemma}
\newtheorem{proposition}[theorem]{Proposition}
\newtheorem{corollary}[theorem]{Corollary}
\theoremstyle{definition}
\newtheorem{definition}[theorem]{Definition}
\newtheorem{remark}[theorem]{Remark}

\newcommand{\N}{\mathbb N}
\newcommand{\Z}{\mathbb Z}
\newcommand{\F}{\mathbb F}
\newcommand{\norm}[1]{\left\lVert #1\right\rVert}
\newcommand{\ip}[2]{\left\langle #1,#2\right\rangle}
\newcommand{\clog}{\operatorname{clog}_2}
\newcommand{\siteH}{H}
\newcommand{\erdosH}{h_{\mathrm E}}
\newcommand{\leanname}[1]{\texttt{#1}}
\hypersetup{
  pdftitle={A candidate N^(1/3+o(1)) upper bound for Erdos Problem 160},
  pdfauthor={Rio Itabe}
}

\title{A candidate \(N^{1/3+o(1)}\) upper bound\\
for Erd\H{o}s Problem 160}
\author{Rio Itabe\\\small Independent researcher}
\date{Preprint, 14 July 2026}

\begin{document}
\maketitle

\begin{center}
\small This manuscript has not been peer reviewed.  It proves an upper-bound
improvement only and does not solve Erd\H{o}s Problem 160.
\end{center}

\begin{abstract}
Let \(\siteH(N)\) be the least number of colours with which
\(\{1,\ldots,N\}\) can be coloured so that every nonconstant four-term
arithmetic progression receives at least three distinct colours.  We give a
concrete construction with explicit palette bounds, after finite choices of a
prime and a field basis, which for every integer \(k\geq1\) and
\[
 N>\max\bigl\{2^{(76k)^2},2^{(40k)^4}\bigr\},
\]
satisfies
\[
 \siteH(N)^{3k}<2{,}985{,}984^kN^{k+3}.
\]
Consequently \(\siteH(N)=O_\varepsilon(N^{1/3+\varepsilon})\) for every
\(\varepsilon>0\), and eventually
\(\siteH(N)\leq N^{1/3+\varepsilon}\).  The construction first uses two
subpower-size digit filters with the property that any progression receiving
at most two colours from their product has the proper pattern \(ABBA\).  If
the two \(ABBA\) equalities also hold in a nine-state carry coordinate, the
associated base-\(p\) digit vectors form a nonconstant affine four-term
progression in \(\F_p^3\).  After an \(\F_p\)-linear identification with
\(\F_{p^3}\), the field norm for the chosen prime \(p>3\) cannot take equal
values simultaneously on the outer pair and on the inner pair.  Hence the
combined carry-and-norm factor is \(ABBA\)-free and has \(9p\) available
colours.  We also prove the exact \(+1\) bridge between this
least-good-colours convention and the
maximum-bad-partitions convention of Erd\H{o}s's original wording.  The stated
theorem chain is formalized in Lean 4.
\end{abstract}

\section{Statement and conventions}

Write \([N]_0=\{0,1,\ldots,N-1\}\), and likewise
\([q]=\{0,1,\ldots,q-1\}\) for a positive integer \(q\).  A nontrivial four-term arithmetic
progression is an ordered quadruple
\[
 (x_0,x_1,x_2,x_3)=(a,a+d,a+2d,a+3d)
 \tag{1.1}
\]
whose entries lie in \([N]_0\) and whose integer step \(d\) is nonzero.  The
translation \(x\mapsto x+1\) identifies this convention with progressions in
\(\{1,\ldots,N\}\).  Allowing the decreasing orientation makes no difference,
since it reverses the same four-element progression.  The Lean definition
\texttt{siteH} uses the corresponding increasing four-element set; translation
and, when necessary, reversal give the same quantified family.

\begin{definition}
For \(N\geq1\), let \(\siteH(N)\) be the least positive integer \(q\) for
which there is a colouring \(c:[N]_0\to[q]\) such that
\[
 \bigl|\{c(x_0),c(x_1),c(x_2),c(x_3)\}\bigr|\geq3
 \tag{1.2}
\]
for every progression \((1.1)\).
\end{definition}

The Lean definition \leanname{siteH} is total on \(\N\); at the unused empty
horizon it has the value \(\leanname{siteH}\ 0=0\).

The maintained Erd\H{o}s Problems page uses this least-good-colours
quantity~\cite{erdosproblems160}.  Erd\H{o}s's original paper instead denotes
by \(h(n)\) the largest number of parts for which every partition is bad:
the union of two parts contains a four-term progression~\cite{erdos1989}.
We denote that original extremum by \(\erdosH(N)\), reserving \(\siteH(N)\)
for the maintained convention.
For \(N\geq4\), the exact finite definition used below is
\[
 \erdosH(N)=
 \max\left\{\,k\in\N:
 2\leq k\leq N\ \text{and every surjective }c:[N]_0\to[k]
 \text{ is bad}\,\right\},
\]
where bad means that some nontrivial four-term progression receives at
most two colours.  Section~9 proves both that the set is nonempty and that
this at-most-two formulation agrees with the literal distinct-nonempty-parts
wording.

\begin{theorem}[Explicit power bound]\label{thm:power}
For integers \(N,k\geq1\), if
\[
 N>\max\bigl\{2^{(76k)^2},2^{(40k)^4}\bigr\},
 \tag{1.3}
\]
then, with \(C=2{,}985{,}984\),
\[
 \siteH(N)^{3k}<C^kN^{k+3}.
 \tag{1.4}
\]
\end{theorem}

\begin{corollary}[One-third upper]\label{cor:epsilon}
For every real \(\varepsilon>0\),
\[
 \siteH(N)=O_\varepsilon\bigl(N^{1/3+\varepsilon}\bigr).
 \tag{1.5}
\]
Moreover, for every \(\varepsilon>0\), eventually
\[
 \siteH(N)\leq N^{1/3+\varepsilon}.
 \tag{1.6}
\]
For \(N\geq4\), the original convention satisfies
\[
 \siteH(N)=\erdosH(N)+1,
 \tag{1.7}
\]
including the literal interpretation in which the two parts are distinct and
every displayed part is nonempty.
\end{corollary}

The maintained page currently records the exponent
\(\log 3/\log 22\approx0.3554\)~\cite{erdosproblems160}.  The family of bounds
in Corollary~\ref{cor:epsilon} has limiting exponent \(1/3\), and Section~8
gives the fixed comparison \(7/20<\log3/\log22\).  The digit-filtering
framework is closely related to Deng--Tidor--Zhao~\cite{dtz} and to Hunter's
explanation of the recorded bound~\cite{huntercomment}.  A literature search
did not locate a prior source for the degree-three norm factor composed with
these filters, but the search was not exhaustive.  We make no claim of
priority or novelty.

\section{The finite pattern reduction}

Suppose four colours use at most two values.  If no three positions have one
common colour, then two distinct colours each occur exactly twice.  Up to
interchanging their names, the only possibilities are
\[
 AABB,\qquad ABAB,\qquad ABBA.
 \tag{2.1}
\]
Thus a label excluding every three-equal pattern and \(ABAB\), together with
an independent label excluding \(AABB\), leaves only the \emph{proper}
pattern
\[
 c_0=c_3,\qquad c_1=c_2,\qquad c_0\ne c_1.
 \tag{2.2}
\]
An additional \(ABBA\)-free label then forces at least three values on the
original progression.  This elementary classification is the formal theorem
\nolinkurl{atMostTwo_noThree_aabb_abab_properABBA}.

We construct three labels:
\begin{enumerate}[label=(\roman*)]
\item one shared mod-24 label \(\lambda_N\), simultaneously excluding
      three-equal patterns and \(ABAB\);
\item a dyadic label \(\mu_N\), excluding \(AABB\);
\item a cubic label \(\nu_N\), excluding \(ABBA\).
\end{enumerate}
Their product is therefore a good colouring.  More explicitly, an
at-most-two tuple of product colours projects to an at-most-two tuple in each
factor.  The first two labels then form a proper-\(ABBA\) filter by the
classification above.  If adjoining the cubic label still left at most two
product colours, the outer and inner equalities forced by that filter would
force the two \(ABBA\) equalities in the cubic coordinate, contradicting its
\(ABBA\)-free property.

\section{Fixed-length digits}

For an integer base \(P>0\), put
\[
 Q_j(n)=\left\lfloor\frac{n}{P^j}\right\rfloor,
 \qquad d_j(n)=Q_j(n)\bmod P.
 \tag{3.1}
\]
If \(0\leq n<P^{L+1}\), then
\[
 v_P(n)=(d_0(n),\ldots,d_L(n))\in\Z^{L+1}
 \tag{3.2}
\]
is a canonical fixed-length representation.  In particular,
\(Q_L(n)<P\), so the top digit is \(d_L(n)=Q_L(n)\).  Equality of all
coordinates in (3.2) therefore implies equality of the represented integers.

The lower digits \(0,\ldots,L-1\) carry auxiliary no-wrap labels.  The top
digit needs no auxiliary label: once the lower quotient induction has been
completed, the progression relation for \(Q_L\) is already the required top
digit relation.  In all three digit constructions below, the square norm is
\[
 \norm{v}^2=\sum_{j=0}^L v_j^2.
 \tag{3.3}
\]

We repeatedly use the same quotient-descent step.  If a no-wrap argument
shows that the current base-\(P\) residues satisfy the progression identities
as integer equalities, substituting \(n_j=r_j+Pq_j\), subtracting those residue
identities, and dividing by \(P\) shows that the quotients \(q_j\) form the
same integer progression.  Iteration exposes every lower digit, and the final
quotient is the top digit described above.

\section{The shared mod-24 filter}

Fix \(B\geq1\) and set
\[
 P=24B+1.
 \tag{4.1}
\]
For \(n<P^{L+1}\), define
\[
 \lambda_{P,L}(n)=
 \left(
   (d_0(n)\bmod24,\ldots,d_{L-1}(n)\bmod24),
   \norm{v_P(n)}^2
 \right).
 \tag{4.2}
\]
The available codomain, not necessarily the attained image, has exactly
\[
 S(P,L)=24^L\bigl((L+1)P^2+1\bigr)
 \tag{4.3}
\]
elements.  The extra endpoint in the norm coordinate is a harmless finite
range containing every possible square norm.

\subsection{A weighted carry-correction lemma}

\begin{lemma}[Weighted digit correction]\label{lem:weighted-carry}
Let \(\alpha,\beta\) be positive integers with
\(\alpha+\beta\leq3\), and suppose
\[
 \alpha x+\beta y=(\alpha+\beta)z.
 \tag{4.4}
\]
Write \(x=r_x+Pq_x\), and similarly for \(y,z\), with
\(0\leq r_x,r_y,r_z<P\).  If
\(r_x\equiv r_y\equiv r_z\pmod {24}\), then
\[
 \alpha r_x+\beta r_y=(\alpha+\beta)r_z
 \tag{4.5}
\]
and
\[
 \alpha q_x+\beta q_y=(\alpha+\beta)q_z.
 \tag{4.6}
\]
\end{lemma}

\begin{proof}
Let
\(E=\alpha r_x+\beta r_y-(\alpha+\beta)r_z\).
Equation (4.4) gives \(P\mid E\), while the residue hypothesis gives
\(24\mid E\).  Since \(P=24B+1\), we have \(\gcd(P,24)=1\), hence
\(24P\mid E\).  On the other hand,
\(\lvert E\rvert<3P<24P\), so \(E=0\).  Subtracting the resulting digit
identity from (4.4) and dividing by \(P\) yields (4.6).
\end{proof}

Iterating Lemma~\ref{lem:weighted-carry} exposes (4.5) in every lower digit
and preserves (4.4) in every quotient.  The top-quotient observation in
Section~3 supplies the last digit.

\subsection{Excluding three equal colours}

Each possible triple of positions in a four-term progression has a positive
weighted relation of the form (4.4):
\[
\begin{array}{c@{\qquad}c@{\qquad}c}
\toprule
\text{positions}&(x,y;z)&(\alpha,\beta)\\
\midrule
012&(x_0,x_2;x_1)&(1,1)\\
123&(x_1,x_3;x_2)&(1,1)\\
013&(x_0,x_3;x_1)&(2,1)\\
023&(x_0,x_3;x_2)&(1,2)\\
\bottomrule
\end{array}
\tag{4.7}
\]
If the three \(\lambda_{P,L}\)-colours agree, Lemma
\ref{lem:weighted-carry} gives digit vectors \(X,Y,Z\) satisfying
\[
 \alpha X+\beta Y=(\alpha+\beta)Z
 \tag{4.8}
\]
coordinatewise, and the norm coordinates give
\(\norm X^2=\norm Y^2=\norm Z^2\).  For every coordinate,
\[
 \alpha(X_j-Z_j)^2+\beta(Y_j-Z_j)^2
 =\alpha X_j^2+\beta Y_j^2-(\alpha+\beta)Z_j^2.
 \tag{4.9}
\]
After summation the right side is zero.  Positivity of the weights and
nonnegativity of the squares force \(X=Y=Z\).  Digit injectivity forces the
three integers to be equal, which is impossible for three distinct positions
of a nontrivial progression.  Thus (4.2) excludes every three-equal pattern.

\subsection{Excluding \(ABAB\) with the same label}

Let \(r_i\) be the current base-\(P\) residues of a four-term progression
and put \(\delta_i=r_{i+1}-r_i\).  The progression relation before taking
residues implies
\[
 \delta_0=\delta_1+uP,\qquad
 \delta_2=\delta_1+vP,\qquad u,v\in\{-1,0,1\}.
 \tag{4.10}
\]
Under an \(ABAB\) equality of (4.2),
\(r_0\equiv r_2\pmod {24}\) and
\(r_1\equiv r_3\pmod {24}\).  Now
\[
 \delta_0-\delta_1=2r_1-r_0-r_2
 \tag{4.11}
\]
is even, whereas \(P\) is odd.  Therefore the first equality in (4.10)
cannot have \(u=\pm1\), and \(u=0\).  The same parity argument applied to
\(\delta_2-\delta_1=r_3+r_1-2r_2\) gives \(v=0\).  Iterating this no-wrap
step through the quotients gives a nonconstant integer vector progression
\[
 X,\quad X+R,\quad X+2R,\quad X+3R,
 \qquad R\ne0.
 \tag{4.12}
\]
The direction \(R\) is nonzero because equality of the fixed-length digit
vectors would, by Section~3, force equality of the represented integers,
contrary to the original nontrivial progression.

Put \(A=\ip XR\) and \(D=\norm R^2\).  The two norm equalities in the
\(ABAB\) pattern give
\[
 \begin{aligned}
 \norm X^2=\norm{X+2R}^2&\quad\Longrightarrow\quad4A+4D=0,\\
 \norm{X+R}^2=\norm{X+3R}^2&\quad\Longrightarrow\quad4A+8D=0.
 \end{aligned}
 \tag{4.13}
\]
Hence \(D=0\), so \(R=0\), a contradiction.  The single label
\(\lambda_{P,L}\) therefore excludes both three-equal patterns and
\(ABAB\); it is counted only once in the final palette.

\section{The dyadic \(AABB\) filter}

This construction uses the separate base \(B\), not the mod-24 base \(P\).
For \(0\leq r<B\), define the bucket
\[
 \beta(r)=\lfloor\log_2(r+1)\rfloor.
 \tag{5.1}
\]
Equality of buckets implies
\[
 y+1<2(x+1),\qquad x+1<2(y+1).
 \tag{5.2}
\]
For \(n<B^{L+1}\), put
\[
 \mu_{B,L}(n)=
 \left(
   (\beta(d_0(n)),\ldots,\beta(d_{L-1}(n))),
   \norm{v_B(n)}^2
 \right).
 \tag{5.3}
\]

\begin{lemma}[Dyadic no-wrap]\label{lem:dyadic}
Suppose four residues \(r_i\in[0,B)\) have congruent consecutive steps
modulo \(B\), and
\[
 \beta(r_0)=\beta(r_1),\qquad
 \beta(r_2)=\beta(r_3).
 \tag{5.4}
\]
Then their three signed representative steps are equal in \(\Z\).
\end{lemma}

\begin{proof}
As in (4.10), write
\(\delta_0=\delta_1+uB\) and
\(\delta_2=\delta_1+vB\), where \(u,v\in\{-1,0,1\}\).
Equivalently,
\[
 2r_1=r_0+r_2+uB,
 \qquad
 r_1+r_3=2r_2+vB.
\]
The inequalities (5.2) give
\[
 r_1\leq2r_0,\qquad r_0\leq2r_1,\qquad
 r_3\leq2r_2,\qquad r_2\leq2r_3.
 \tag{5.5}
\]
If \(u=-1\), then \(B=r_0+r_2-2r_1\); since \(B>r_2\), this gives
\(r_0>2r_1\), contradicting (5.5).  Suppose \(u=1\).  Since \(B>r_1\),
the first displayed identity gives \(r_1>r_0+r_2\).  If \(v=-1\), then
\(B=2r_2-r_1-r_3>r_2\), so \(r_2>r_1+r_3\), contradicting
\(r_1>r_2\).  If \(v=0\), then \(r_3=2r_2-r_1\geq0\), whence
\(r_2>r_0\) and therefore \(r_1>2r_0\), contradicting (5.5).  If
\(v=1\), the second displayed identity and \(B>r_1\) give
\(r_3>2r_2\), again contradicting (5.5).  Thus \(u=0\).

Now \(2r_1=r_0+r_2\).  If \(v=-1\), then
\(r_3=2r_2-r_1-B<r_2-r_1\).  Together with \(r_2\leq2r_3\), this
gives
\[
 2r_1<r_2\leq r_0+r_2=2r_1,
\]
a contradiction.  If \(v=1\), then
\(B=r_1+r_3-2r_2\).  The inequalities \(B>r_3\) and
\(B>r_0=2r_1-r_2\) give, respectively,
\(r_1>2r_2\) and \(r_3>r_1+r_2\).  Since \(r_3\leq2r_2\), the latter
implies \(r_1<r_2\), a contradiction.  Hence \(v=0\).
\end{proof}

Quotient induction and the top digit again turn an \(AABB\) equality of
\(\mu_{B,L}\) into an integer vector progression (4.12).  Its direction is
nonzero by the same fixed-length digit injectivity argument.  For
\(Q(t)=\norm{X+tR}^2\), the two adjacent-pair norm equalities give
\[
 Q(0)=Q(1)\Longrightarrow2A+D=0,\qquad
 Q(2)=Q(3)\Longrightarrow2A+5D=0.
 \tag{5.6}
\]
Thus \(D=0\), contradicting \(R\ne0\).  Hence (5.3) is \(AABB\)-free.

The available codomain has cardinality
\[
 D(B,L)=\bigl(\lfloor\log_2B\rfloor+1\bigr)^L
         \bigl((L+1)B^2+1\bigr).
 \tag{5.7}
\]

\section{All-horizon parameters and subpower bounds}

Let
\[
 q=\clog N,\qquad
 m=\lfloor\sqrt{q-1}\rfloor+1,\qquad
 L=m-1,\qquad B=2^m,\qquad P=24B+1.
 \tag{6.1}
\]
Here \(\clog N\) is the least \(q\) with \(N\leq2^q\).  The elementary
inequality \(q\leq m^2\) yields
\[
 N\leq2^{m^2}=B^{L+1}\leq P^{L+1},
 \tag{6.2}
\]
so both preceding digit labels cover \([N]_0\).

For the shared mod-24 factor, (4.3) becomes
\[
 S_N=24^{m-1}\bigl(m(24\cdot2^m+1)^2+1\bigr)
     \leq2^{8m+11}.
 \tag{6.3}
\]
Indeed, \(24^{m-1}\leq2^{5m}\),
\(P\leq2^{m+5}\), \(m\leq2^m\), and therefore
\(mP^2+1\leq2^{3m+11}\).

For the dyadic factor, (5.7) gives the exact available cardinality
\[
 D_N=(m+1)^{m-1}\bigl(m2^{2m}+1\bigr).
 \tag{6.4}
\]
With \(r=\lceil\log_2(m+1)\rceil\), elementary estimates give
\[
 D_N\leq2^{mr+2m+1}\leq2^{5m\sqrt m}.
 \tag{6.5}
\]

We use the following explicit thresholds.  A generic palette bounded by
\(2^{Am+B_0}\) has \(K\)-th power
less than \(N\) once
\[
 N>2^{T^2},\qquad T=\max\{1,2(A+B_0)K\}.
 \tag{6.6}
\]
To see this, put \(s=\lfloor\sqrt{q-1}\rfloor\).  The threshold gives
\(s\geq T\), while \(m=s+1\leq2s\), so
\((Am+B_0)K\leq s^2\leq q-1\).  Applying (6.6) to (6.3) with
\((A,B_0,K)=(8,11,2k)\) gives
\[
 N>2^{(76k)^2}\quad\Longrightarrow\quad S_N^{2k}<N.
 \tag{6.7}
\]

For (6.5), the sufficient elementary estimate used in the Lean formalization is
\[
 N>2^{(20K)^4}\quad\Longrightarrow\quad D_N^K<N.
 \tag{6.8}
\]
It follows by taking \(s=\lfloor\sqrt{q-1}\rfloor\), observing that
\(s\geq(20K)^2\), and using
\(5(s+1)\sqrt{s+1}\,K\leq s^2\leq q-1\).
With \(K=2k\),
\[
 N>2^{(40k)^4}\quad\Longrightarrow\quad D_N^{2k}<N.
 \tag{6.9}
\]
Since \(S_N,D_N>0\), (6.7) and (6.9) imply
\[
 (S_ND_N)^k<N;
 \tag{6.10}
\]
one may square both sides and use
\(((S_ND_N)^k)^2=S_N^{2k}D_N^{2k}<N^2\).

\section{The cubic proper-\(ABBA\) factor}

The preceding product leaves only the proper pattern (2.2).  We now exclude
it with \(O(N^{1/3})\) available colours.

\subsection{A nine-state carry shield}

Let \(p>3\) be prime.  For a residue \(0\leq r<p\), define its third
\[
 \tau_p(r)=
 \begin{cases}
 0,&3r<p,\\
 1,&p\leq3r<2p,\\
 2,&2p\leq3r.
 \end{cases}
 \tag{7.1}
\]
Write every \(n<p^3\) uniquely as
\[
 n=r+p(w+pu),\qquad0\leq r,w,u<p,
 \tag{7.2}
\]
and put
\[
 \sigma_p(n)=(\tau_p(r),\tau_p(w))\in\{0,1,2\}^2.
 \tag{7.3}
\]
Only the two carry-bearing digits are labelled, so \(\sigma_p\) has nine
available values.

\begin{lemma}[Thirds no-wrap]\label{lem:thirds}
Let \(r_0,r_1,r_2,r_3\in[0,p)\), with their consecutive signed differences
congruent modulo \(p\).  If
\[
 \tau_p(r_0)=\tau_p(r_3),\qquad
 \tau_p(r_1)=\tau_p(r_2),
 \tag{7.4}
\]
then the three representative differences are equal in \(\Z\).
\end{lemma}

\begin{proof}
Let the common outer region in (7.4) be \(a\), and the common inner region
be \(b\).  Write
\[
 3r_j=\rho_jp+e_j,
 \qquad 0\leq e_j<p,
 \qquad
 \rho_0=\rho_3=a,\qquad \rho_1=\rho_2=b.
\]
Every representative difference lies strictly between \(-p\) and \(p\).
Consequently, relative to the middle difference, there are
\(u,v\in\{-1,0,1\}\) such that
\[
 \delta_0=\delta_1+up,
 \qquad
 \delta_2=\delta_1+vp.
\]
Put \(d=b-a\in\{-2,-1,0,1,2\}\).  Multiplying the two difference
identities by \(3\) and using the displayed decompositions gives
\[
 (d-3u)p=e_0+e_2-2e_1,
 \qquad
 (-d-3v)p=2e_2-e_1-e_3.
 \tag{7.5}
\]
Both right-hand sides lie strictly between \(-2p\) and \(2p\).  Since the
coefficients on the left are integers, the only possibilities are
\[
\begin{array}{c@{\qquad}cc}
\toprule
d&u&v\\
\midrule
-2&-1&1\\
-1&0&0\\
0&0&0\\
1&0&0\\
2&1&-1\\
\bottomrule
\end{array}
\]
When \(d=-2\), subtracting the second equation in (7.5) from the first
would give
\[
 2p=e_0+e_3-e_1-e_2<2p,
\]
which is impossible.  When \(d=2\), subtracting the first equation from the
second similarly gives
\[
 2p=e_1+e_2-e_0-e_3<2p.
\]
Thus the two extreme rows are impossible, and all remaining rows have
\(u=v=0\), as required.  This argument is formalized as
\leanname{thirdsRegion\_noWrap}.
\end{proof}

Apply Lemma~\ref{lem:thirds} first to the low digits in (7.2).  To spell out
the quotient descent, write \(n_j=r_j+pq_j\).  The original progression
identities and the exact residue identities give
\[
 p(q_0+q_2-2q_1)=0,
 \qquad
 p(q_1+q_3-2q_2)=0.
\]
Since \(p\ne0\), the quotients \(q_j=s_j+pu_j\) form an integer
progression.  Apply the lemma again to the middle digits; the same division
then leaves an integer progression in the high digits.  Thus, whenever a
nontrivial four-term progression has \(ABBA\) equality under the shield
(7.3), its integer digit vectors form an affine progression.  Reducing the
three coordinates modulo \(p\) gives
\[
 X,\quad X+R,\quad X+2R,\quad X+3R
 \tag{7.6}
\]
over \(\F_p^3\), with \(R\ne0\).  The shield does not make arbitrary
progressions carry-free; this conclusion specifically uses the two \(ABBA\)
shield equalities.  For completeness, \(R\) remains nonzero after reduction:
the integer digit step is nonzero by injectivity of the fixed-length
representation, every coordinate has absolute value less than \(p\), and so
coordinatewise divisibility by \(p\) would force that step to be zero.

\subsection{The degree-three norm}

Let
\[
 \overline D_p(n)=(r\bmod p,w\bmod p,u\bmod p)\in\F_p^3
\]
be the reduced digit vector from (7.2).  Choose an \(\F_p\)-linear
isomorphism
\(\iota_p:\F_p^3\xrightarrow{\sim}\F_{p^3}\), and let
\[
 \mathcal N:\F_{p^3}\longrightarrow\F_p
 \tag{7.7}
\]
be the field norm.  On every affine line it is a cubic polynomial:
\[
 \mathcal N(X+tR)=at^3+bt^2+ct+d,\qquad a=\mathcal N(R).
 \tag{7.8}
\]
This follows, for example, by writing the norm as the determinant of the
\(3\times3\) matrix of multiplication and taking the leading determinant.

Apply \(\iota_p\) to (7.6), and denote the resulting base point and nonzero
direction again by \(X,R\).  Suppose their four norm values have the
\(ABBA\) pattern.  From \(f(0)=f(3)\) and \(f(1)=f(2)\), where the right
side of (7.8) is \(f(t)\),
we obtain
\[
 27a+9b+3c=0,\qquad7a+3b+c=0.
 \tag{7.9}
\]
Subtracting three times the second equation from the first gives \(6a=0\).
Because \(p>3\), the element \(6\) is nonzero in \(\F_p\), so
\(\mathcal N(R)=a=0\).  A field norm vanishes only at zero, contradicting
\(R\ne0\).

It follows that
\[
 \nu_p(n)=
 \left(\sigma_p(n),
 \mathcal N\bigl(\iota_p(\overline D_p(n))\bigr)\right)
 \tag{7.10}
\]
is \(ABBA\)-free on \([0,p^3)\), with exactly \(9p\) available codomain
values.

\subsection{A prime for every horizon and the constant}

Let \(q=\clog N\), and define
\[
 e=\left\lfloor\frac{q+2}{3}\right\rfloor+2,\qquad y=2^e.
 \tag{7.11}
\]
Then \(y>3\) and \(N\leq y^3\).  Bertrand's postulate supplies a prime
\(p\) with
\[
 y<p\leq2y,
 \tag{7.12}
\]
so \(N\leq p^3\) and (7.10) covers \([N]_0\).  If \(N>1\), then
\(3e\leq q+8\) and \(2^{q-1}<N\).  Consequently
\[
 y^3\leq2^{q+8}<512N,\qquad p^3<4096N.
 \tag{7.13}
\]
Therefore the cubic available palette \(U_N=9p\) satisfies
\[
 U_N^3=(9p)^3<729\cdot4096\,N
 =2{,}985{,}984N.
 \tag{7.14}
\]
No optimization of this uniform constant is intended.

\section{Proof of the main bound}

Take the product colouring
\[
 c_N(n)=\bigl(\lambda_N(n),\mu_N(n),\nu_N(n)\bigr).
 \tag{8.1}
\]
The shared label excludes three-equal patterns and \(ABAB\), the dyadic
label excludes \(AABB\), and the cubic label excludes the remaining proper
\(ABBA\).  Section~2 therefore shows that every nontrivial four-term
progression receives at least three product colours.  Hence
\[
 \siteH(N)\leq S_ND_NU_N.
 \tag{8.2}
\]

Assume (1.3).  Equations (6.10) and (7.14) give
\[
\begin{aligned}
 \siteH(N)^{3k}
 &\leq(S_ND_N)^{3k}U_N^{3k}\\
 &<N^3\bigl(2{,}985{,}984N\bigr)^k\\
 &=2{,}985{,}984^kN^{k+3},
\end{aligned}
\tag{8.3}
\]
which proves Theorem~\ref{thm:power}.

Taking positive \(3k\)-th roots gives
\[
 \siteH(N)<C^{1/3}N^{1/3+1/k}.
 \tag{8.4}
\]
Given \(\varepsilon>0\), choose \(k\) with \(1/k<\varepsilon\).  Equation
(8.4) first yields (1.5).  For (1.6), put
\(\delta=\varepsilon-1/k>0\); eventually
\(C^{1/3}\leq N^\delta\), so the fixed coefficient is absorbed into the
remaining exponent.

As a completely rational visible comparison, \(k=60\) gives exponent
\((60+3)/(180)=7/20\).  The integer inequality
\[
 22^7<3^{20}
 \tag{8.5}
\]
implies
\(7/20<\log3/\log22\), so even this fixed specialization improves the
exponent recorded on the maintained page.

\section{The original partition convention}

We now justify (1.7) rather than silently identifying two different
definitions.  Fix \(N\geq4\) and write \(H=\siteH(N)\).  First,
\(H\leq N\): the injective \(N\)-colouring is good because every nontrivial
progression has four distinct entries.  Choose a good \(H\)-palette
colouring and compress its codomain to its image, of size \(\ell\leq H\).
The compressed colouring is surjective and good.  Minimality of \(H\) forces
\(\ell=H\), so a surjective good \(H\)-colouring exists.

If a surjective good colouring has \(k<N\) colours, one fibre has at least
two elements.  Splitting one element into a fresh colour preserves goodness.
Iteration therefore gives a surjective good colouring for every
\(H\leq k\leq N\).  Conversely, if \(k<H\), minimality of \(H\) says that
no \(k\)-colouring is good, so in particular every surjective
\(k\)-colouring is bad.

No two-colouring is good because the progression \(0,1,2,3\) itself uses at
most two colours.  Thus \(H\geq3\), the finite set defining \(\erdosH(N)\)
is nonempty, and the preceding two paragraphs show that its maximum is
exactly \(H-1\geq2\).

Finally, when \(2\leq k\leq N\), the literal phrase ``the union of two
distinct nonempty parts contains a progression'' is equivalent to saying
that some progression uses at most two colours.  If it uses two colours,
choose those two parts.  If it uses only one, choose that part and any other
nonempty part.  The converse is immediate.  This proves
\[
 \erdosH(N)=\siteH(N)-1\qquad(N\geq4),
 \tag{9.1}
\]
under both readings, and hence (1.7).  The corresponding original-convention
version of (1.4) is
therefore
\[
 (\erdosH(N)+1)^{3k}<C^kN^{k+3}.
 \tag{9.2}
\]

\section{Lean formalization}

The Lean development uses Lean \(4.30.0\) and mathlib \(4.30.0\), pinned by
the repository manifest.  The principal theorem names are listed below.
\[
\begin{array}{p{0.36\textwidth}p{0.56\textwidth}}
\toprule
\text{Mathematical role}&\text{Lean theorem}\\
\midrule
two-colour classification&
\nolinkurl{atMostTwo_noThree_aabb_abab_properABBA}\\
shared mod-24 filter&
\nolinkurl{horizonNoThreeMod24Colour_isNoThreeEqualColouring},
\nolinkurl{horizonABABMod24Colour_isABABFreeColouring}\\
dyadic \(AABB\) filter&
\nolinkurl{horizonDyadicSquareColour_isAABBFreeColouring}\\
cubic \(ABBA\) factor&
\nolinkurl{horizonCubicColour_abba_free}\\
displayed threshold bridge&
\nolinkurl{explicitUnifiedOneThirdThreshold_eq_displayed}\\
exact natural-power bound&
\nolinkurl{siteH_oneThirdPlusSubpower_unified_displayed}\\
real epsilon Big-O&
\nolinkurl{siteH_oneThirdPlusEpsilon_unified_isBigO}\\
coefficient-one eventual bound&
\nolinkurl{siteH_le_rpow_oneThird_add_epsilon_eventually}\\
strict source bridge&
\nolinkurl{siteH_eq_sourceSurjectiveOriginalH_add_one}\\
literal distinct-two characterization&
\nolinkurl{sourceSurjectiveOriginalH_literal_distinct_two_characterization}\\
\bottomrule
\end{array}
\tag{10.1}
\]

The repository check builds the Lean project, scans the E160 sources for
\leanname{sorry} and \leanname{admit}, prints the axioms of the listed
theorems, and verifies the hashes of short source excerpts.  The reported
axioms are standard Lean/mathlib principles such as propositional
extensionality, classical choice, and quotient soundness; no project-defined
axiom or \leanname{sorryAx} is used.

Kernel checking establishes the formal implications represented by the code;
source interpretation, novelty, and peer review lie outside that check.  No
asymptotic colouring theorem from~\cite{dtz} is imported: the mod-24 and
dyadic constructions used above are proved explicitly in the Lean development.

\section{LLM assistance}

LLMs accessed through ChatGPT Pro and OpenAI Codex, primarily GPT-5.6, were
used during proof exploration, Lean formalization, testing, and editing.  The
author is responsible for the manuscript and code; LLM outputs were not
treated as proof or independent review.

\begin{thebibliography}{9}

\bibitem{erdos1989}
P.~Erd\H{o}s,
\emph{Some problems and results on combinatorial number theory},
in Graph Theory and Its Applications: East and West, Annals of the New York
Academy of Sciences 576 (1989), 132--145.  Problem passage on printed
p.~143.

\bibitem{dtz}
M.~Deng, J.~Tidor, and Y.~Zhao,
\emph{Uniform sets with few progressions via colorings},
Mathematical Proceedings of the Cambridge Philosophical Society 179 (2025),
79--103; \href{https://arxiv.org/abs/2307.06914}{arXiv:2307.06914}.

\bibitem{erdosproblems160}
T.~Bloom, ed.,
\emph{Erd\H{o}s Problem 160}, maintained problem page,
\url{https://www.erdosproblems.com/160}, accessed 14 July 2026.

\bibitem{huntercomment}
Z.~Hunter,
comment on Erd\H{o}s Problem 160,
\url{https://www.erdosproblems.com/forum/thread/160}, posted 18 October 2025
(displayed time has no stated timezone), accessed 14 July 2026.

\bibitem{lean}
L.~de Moura and S.~Ullrich,
\emph{The Lean 4 theorem prover and programming language},
Automated Deduction -- CADE 28, Lecture Notes in Computer Science 12699
(2021), 625--635.

\end{thebibliography}

\end{document}
